\section{La definición mínima de Agent: entidad, entorno y actividad orientada a metas}

La definición de Agent que da Su Yu (苏煜) es muy sencilla: un Agent es una entidad con límites definidos que opera en un entorno externo y realiza goal-directed activities, es decir, actividades orientadas a metas. La definición es amplia a propósito, porque animales, personas, robots, agentes que operan páginas web y coding agents caben en este marco. Lo decisivo no es si se parece a un humano, sino si tiene límites identificables, entradas del entorno y un comportamiento impulsado por metas.

\begin{figure}[H]
\centering
\includegraphics[width=0.92\textwidth]{agent-loop.png}
\caption{El ciclo mínimo de un Agent: entorno, observación, razonamiento, acción y memoria.}
\end{figure}

\begin{knowledgebox}{Lectura de la figura: cómo leer el ciclo del Agent}
En el extremo izquierdo de la figura está el entorno, que puede ser una página web, un escritorio, un repositorio de código, una base de datos o el mundo de un robot; arriba, Observation es el estado percibido; en el centro, el Agent razona a partir de la meta, la observación y la memoria; a la derecha, Action modifica el entorno; abajo, Memory/World Model consolida la experiencia de interacción en un estado reutilizable. La conclusión que sustenta es esta: un Agent no es una respuesta única, sino un sistema de lazo cerrado.
\end{knowledgebox}

\begin{importantbox}{Una fórmula simplificada}
Una acción de un Agent puede escribirse como:
$$
a_t = \pi(o_t, m_t, g)
$$
donde $a_t$ es la acción del paso $t$; $o_t$ es la observación actual; $m_t$ es la memoria o el modelo del mundo; $g$ es la meta; $\pi$ es la política o función de decisión. Esta fórmula no aparece explícitamente en la entrevista; es una condensación didáctica de la definición de Su Yu.
\end{importantbox}

\subsection{Por qué la definición debe ser amplia}

Si la definición es demasiado estrecha, es fácil confundir un Agent con una forma de producto concreta, como un agente web, un agente de escritorio o un coding agent. La definición de Su Yu considera todos estos casos como instancias de una misma clase de sistema en entornos distintos. Browser use, desktop use, mobile use, GUI, CLI, API y coding son solo means to an end; la meta real es un universal digital agent capaz de completar todo tipo de tareas en el digital world.

\begin{warningbox}{Malentendido frecuente: un Agent no es una ventana de chat con plugins}
Que una ventana de chat invoque una herramienta es apenas un corte temprano de lo que es un Agent. Un Agent real debe percibir el estado, planificar acciones, ejecutarlas, aprovechar la retroalimentación y mantener la meta a lo largo de tareas de varios pasos. Reducir un Agent a «un LLM que sabe invocar APIs» subestima la importancia de memory, world model, reliability y cost.
\end{warningbox}

\subsection{Resumen de la sección}

La definición mínima de Agent es: una entidad con límites definidos que actúa en un entorno con base en metas. Esta definición permite colocar en una misma línea histórica los primeros sistemas simbólicos, el aprendizaje por refuerzo, el semantic parsing, el uso de herramientas por LLM y los coding agents modernos.
